\documentclass{article}
\usepackage{amsmath,amssymb}
\usepackage{xurl}
\usepackage[hidelinks]{hyperref}
\setlength{\emergencystretch}{2em}
% Shared commands for notes in docs/paper-gaps/.

\DeclareUnicodeCharacter{2080}{\ifmmode{}_{0}\else\textsubscript{0}\fi}
\DeclareUnicodeCharacter{2081}{\ifmmode{}_{1}\else\textsubscript{1}\fi}
\DeclareUnicodeCharacter{2082}{\ifmmode{}_{2}\else\textsubscript{2}\fi}
\DeclareUnicodeCharacter{2099}{\ifmmode{}_{n}\else\textsubscript{n}\fi}

\newcommand{\Lean}{\textsc{Lean}}
\ExplSyntaxOn
\NewDocumentCommand{\leanid}{m}{
  \ifmmode
    \textsf{\detokenize{#1}}
  \else
    \group_begin:
    \urlstyle{sf}
    \tl_set:Nn \l_tmpa_tl {#1}
    \tl_replace_all:Nnn \l_tmpa_tl {\_} {_}
    \exp_args:NV \path \l_tmpa_tl
    \group_end:
  \fi
}
\ExplSyntaxOff
\providecommand{\tr}{\operatorname{tr}}
\newcommand{\tnleanIssueBase}{https://github.com/LionSR/TNLean/issues/}
\newcommand{\tnleanPrBase}{https://github.com/LionSR/TNLean/pull/}
\newcommand{\tnleanDocBase}{https://sirui-lu.com/TNLean/docs/}
\newcommand{\ghissue}[1]{\href{\tnleanIssueBase#1}{GitHub issue \##1}}
\newcommand{\ghpr}[1]{\href{\tnleanPrBase#1}{PR \##1}}
\newcommand{\doclink}[1]{\href{\tnleanDocBase#1.html}{\path{#1}}}

% Dirac notation and the identity operator, as in blueprint/src/macros/common.tex.
% \providecommand keeps note-local definitions authoritative.
\usepackage{dsfont}
\providecommand{\ket}[1]{|#1\rangle}
\providecommand{\bra}[1]{\langle #1|}
\providecommand{\braket}[2]{\langle #1 | #2 \rangle}
\providecommand{\ketbra}[2]{|#1\rangle\!\langle #2|}
\providecommand{\Id}{\mathds{1}}
\providecommand{\one}{\mathds{1}}
\providecommand{\C}{\mathbb{C}}
\providecommand{\MN}[1]{M_{#1}(\C)}
\providecommand{\MD}{M_D(\C)}

% Verdict marker: \gapnote{<kind>}{<status>}. Kind is the policy
% classification (clarification, local-correction, scope-restriction,
% unfaithful, false-source, open-gap); status is open, wip (elimination
% underway), resolved, or historical. Machine-read by the site index;
% prints a verdict line.
\newcommand{\gapnote}[2]{\par\noindent\textbf{Verdict:} #1 (#2).\par}

\title{Half-Chain Spectrum: Canonical Orientation, Normalization, and Decay}
\author{}
\date{2026-10-04}
\begin{document}
\maketitle
\gapnote{scope-restriction}{open}

\section{Source and index convention}
Theorem~6 of~\cite{PerezGarcia2007MPSReps}, preceded by the calculation
at lines 970--985 of \path{Papers/quant-ph_0608197/MPSarchive.tex}, identifies
the limiting half-chain reduced spectrum with that of
$\Lambda\otimes\Lambda$. The canonical convention is
$\mathcal E_A(I)=I$ and $\mathcal E_A^*(\Lambda)=\Lambda$, where
$\Lambda=\operatorname{diag}(\lambda_\alpha)>0$ and
$\operatorname{tr}\Lambda=1$.

For $\mathcal E_A^L(X)\to\operatorname{tr}(\Lambda X)I$, the source's
matrix element gives
\begin{align}
    \langle\alpha'|\mathcal E_A^L(|\beta'\rangle\langle\beta|)|\alpha\rangle
     & \longrightarrow
    \lambda_\beta\delta_{\alpha\alpha'}\delta_{\beta\beta'}.
    \label{eq:halfchain_note_unital_gram}
\end{align}
The printed display at line 979 instead places the weight on $\alpha$.
That placement is correct for the trace-preserving convention
$\mathcal E_A^L(X)\to\operatorname{tr}(X)\Lambda$.
Taking adjoints of every Kraus matrix exchanges the conventions and swaps
and transposes the Gram matrix. In both cases the limiting two-boundary
product is $\lambda_\alpha\lambda_\beta$.
The final unital theorem concerns the original tensor and its actual ring
state; it does not replace the physical state by the adjoint-family state.

\section{Finite-size normalization and zero eigenvalues}
Let $\Psi_L$ be the half-chain boundary family, $G_L=\Psi_L^\dagger\Psi_L$,
and $W_L=SG_L^{\mathsf T}S$. The physical unnormalized reduced density is
$\rho_L=\Psi_LW_L\Psi_L^\dagger$. The positive representative
$H_L=\sqrt{G_L}W_L\sqrt{G_L}$ satisfies
\begin{align}
    X^{D^2}\chi_{\rho_L}(X) & =X^{d^L}\chi_{H_L}(X),
    \label{eq:halfchain_note_charpoly}               \\
    Z_L:=\operatorname{tr}\rho_L
                            & =\operatorname{tr}H_L
    =\sum_{\sigma,\tau}|V^{(2L)}(A)_{\sigma\tau}|^2.
    \label{eq:halfchain_note_norm}
\end{align}
Normalize both matrices by this actual $Z_L$, rather than assuming it is
one at finite length. Convergence $H_L\to\Lambda\otimes\Lambda$ gives
$Z_L\to1$, hence eventual strict positivity. At a zero-norm length the
normalized matrices are defined as zero.

List eigenvalues in decreasing order with multiplicity, then append zeros.
Identity~(\ref{eq:halfchain_note_charpoly}) identifies the two resulting
sequences at every finite length, including $d^L<D^2$. The physical list
is zero beyond its dimension and beyond $D^2$.

\section{Decay statement}
The source states corrections of order $|\nu_2|^L$. Without semisimplicity
at the subleading eigenvalues, powers of the complementary transfer map
can carry polynomial factors from Jordan blocks. The geometric estimate
used here chooses a radius strictly above the complementary spectral
radius and below one, and bounds all Gram entries by $Cr^L$.
This suffices for the stated spectral convergence.

Continuity of the positive square root applies on the entire positive
cone, including singular limits. Weyl's inequality then gives convergence
of each ordered eigenvalue of $H_L/Z_L$; the exact characteristic-polynomial
identity transfers it to $\rho_L/Z_L$. No sharp quantitative rate for the
square-root step, entropy convergence, or infinite-dimensional reduced
operator is asserted.

\section{Formal statements}
The finite-size comparison and normalized spectral convergence are
\leanid{MPSTensor.halfChainEigenvalues_eq_boundary} and
\leanid{MPSTensor.tendsto_halfChainEigenvalues}.
The normal-tensor unital statement is
\leanid{MPSTensor.tendsto_halfChainEigenvalues_unital}; it assumes
normality explicitly. Bare uniqueness of the peripheral eigenvalue does
not imply irreducibility: the identity channel has peripheral spectrum
$\{1\}$ but a non-scalar fixed-point space.

The source context includes all three one-block canonical conditions from
Theorem~4, lines 742--759: unitality, a faithful dual fixed point, and
scalar fixed points. Together with Condition C2 and its spectral-radius-one
normalization at lines 954--963, these imply normality by
\leanid{MPSTensor.isNormalTensor_of_pgvwc07_canonical_c2}.
The explicit source-hypothesis conclusion is
\leanid{MPSTensor.tendsto_halfChainEigenvalues_of_pgvwc07_canonical_c2}.
For an explicitly one-block canonical tensor, no irreducibility or
complementary spectral gap is assumed in addition to the three canonical
conditions and normalized C2 hypothesis.
The earlier normal-tensor result is a valid intermediate theorem, not by
itself the formal statement from those source hypotheses.

The more general explicit-gap theorem allows singular nonnegative fixed
points. Such points are outside the faithful canonical hypotheses of the
source and are not asserted to be normal.

\section{Remaining weighted multiblock scope}
The canonical form in Theorem~4 is a weighted direct sum; the three
canonical conditions apply separately to its blocks. Condition C2 does
not alone make the original representation a single faithful unital block.
For example, let $0<r<1$ and
$A_0=\operatorname{diag}(1,0)$, $A_1=\operatorname{diag}(0,r)$.
The transfer spectrum is $\{1,r^2,0,0\}$, with a simple peripheral
eigenvalue one, but the canonical form has two weighted blocks and the
ring vector is $|0^N\rangle+r^N|1^N\rangle$.
The phrase ``just one block'' at source lines 960--963 therefore requires
an asymptotic dominant-block interpretation rather than literal deletion
of all subleading weighted blocks at finite length.

The new source-hypothesis corollary treats an explicitly one-block
canonical tensor. It repairs the previously missing derivation of
normality in that setting; it does not formalize the reduction from an
arbitrary weighted canonical representation satisfying C2 to its dominant
block. That reduction, including the disappearing subleading physical
contribution and spectral normalization, remains open here. Accordingly,
the unrestricted source theorem is not marked as completely formalized.
The boundary-index and decay corrections in the preceding sections remain
resolved independently of this representation-scope obligation.

\bibliographystyle{alpha}
\bibliography{references}
\end{document}
